\section{Design Science Research}
\label{sec:design-science-research}

Information systems research that produces a working artefact faces a structural question: is building the system itself research, or only the activity that surrounds it. Design Science Research (DSR) answers that the artefact and the knowledge generated by constructing it are inseparable, since ``knowledge and understanding of a design problem and its solution are acquired in the building and application of an artifact'' \parencite[p.~82]{hevner2004design}. \textcite{hevner2004design} frames DSR as the paradigm in which utility, rather than truth, is the goal, and which must produce a viable construct, model, method, or instantiation. \textcite{peffers2007dsrm} operationalise this paradigm as a six-activity process: problem identification, definition of objectives, design and development, demonstration, evaluation, and communication. The two contributions are complementary. Hevner sets out what DSR is and how its quality should be judged; Peffers et al. specify how to carry out and report a DSR project.

DSR fits this thesis because the contribution is dual. The artefact is Ressursplanlegger, a web-based planning platform that produces daily driver and vehicle assignments for traffic coordinators in Norwegian transport companies. The knowledge produced alongside it concerns how algorithm-assisted planning can be structured for small and medium-sized transport firms whose current practice is manual and dependent on coordinator memory. A pure case study would describe the practice without changing it. Pure software development would build a system without producing transferable claims. DSR is chosen because both outputs are needed: a system that can be inspected and a set of design choices, constraint formulations, and benchmarking results that survive beyond this specific implementation.

A further methodological choice concerns how the artefact is assessed. \textcite{wieringa2014dsm} draws a sharp line between validation and evaluation: ``To validate a treatment is to justify that it would contribute to stakeholder goals if implemented [...] Validation is contrasted with evaluation, which is the investigation of a treatment as applied by stakeholders in the field'' \parencite[p.~31]{wieringa2014dsm}. The decisive criterion is researcher control. In validation the researchers experiment with a model of how the artefact would be used; in evaluation, stakeholders use the artefact independently. Ressursplanlegger has not been deployed at a transport company, so coordinators have not used it in their own daily work. This thesis therefore validates the artefact through algorithm benchmarking on synthetic datasets and through requirements traceability against the needs identified in interviews. It does not evaluate it. The consequence is that the thesis can claim the system would address the documented needs and meets its functional requirements, but cannot claim observed effects on real planning work; that claim would require a field deployment.

The role of this section is to establish DSR as the methodological frame for the project. \Cref{chap:method} applies the six activities of \textcite{peffers2007dsrm} to the concrete steps of this thesis, specifies how interviews, requirements work, system development, and benchmarking instantiate the design cycle, and discusses the validity threats that follow from the validation rather than evaluation choice.